\newpage
\addcontentsline{toc}{chapter}{Abstract}
\chapter*{Abstract\markboth{Abstract}{}}
Household food waste is one of the most pressing socio-economic and environmental challenges of the modern era, with domestic kitchens accounting for over 60\% of global edible food wastage. A primary driver is culinary friction—home cooks frequently discard usable ingredients because they struggle to identify what dishes can be prepared from heterogeneous, partially depleted pantry items before spoilage occurs.

To address this challenge, this project presents \textbf{PlateWise}, an individual full-stack AI-powered culinary intelligence and zero-waste kitchen management platform. PlateWise operationalizes computer vision and large language models (LLMs) into an intuitive, high-performance web application. Users capture or drag-and-drop unstructured photos of refrigerator shelves, kitchen countertops, or grocery hauls. The system processes the image through client-side HTML5 Canvas downscaling and sends base64 payloads to a serverless edge backend on Vercel. A multi-modal visual inference cascade powered by Groq's low-latency LPU cloud (utilizing \texttt{qwen/qwen3.8-27b} with candidate fallbacks to \texttt{meta-llama/llama-4-scout} and \texttt{llama-3.2-11b-vision}) detects up to 16+ distinct produce and pantry ingredients simultaneously in under 1.5 seconds.

The detected inventory is fed into a context-aware culinary reasoning engine (\texttt{openai/gpt-oss-120b} / \texttt{llama-3.3-70b-versatile}) that dynamically factors in time of day (Breakfast, Lunch, Dinner), seasonal availability (Monsoon, Summer, Winter), and regional cuisine styles (Indian Comfort, Mediterranean, Continental) to generate four structured gourmet recipes. The platform further incorporates an automated Pantry Audit Engine that computes mathematical compatibility percentages, identifies missing items, and provides intelligent culinary substitutions. An interactive Cooking Workstation embeds real-time YouTube Masterclass video tutorials and interactive progress checklists, while an automated Recipe Vault mirrors completed meals to client \texttt{localStorage} to ensure complete resilience against serverless cold-start data loss.

Comprehensive empirical evaluation demonstrates sub-2-second vision detection latency, robust handling of real-world multi-ingredient scenes, and seamless full-pipeline execution from visual capture to cooked recipe archiving. The production application is publicly deployed on Vercel at \url{https://platewise1705.vercel.app/} with open-source code on GitHub.
\vspace{6mm}
\hrule